% !TeX program = xelatex
% UTF-8 transcription; source wording and values are retained without correction.
% Google Drive: Notebook/HTN; originals in images/ (unchanged JPEG bytes).
% pic01.jpg = 20261012_191556507.jpg
% pic02.jpg = 20261012_191603503.jpg
% pic03.jpg = 20261012_191608346.jpg
% pic04.jpg = 20261012_191616426.jpg
% pic05.jpg = 20261012_191622309.jpg
% pic06.jpg = 20261012_191629434.jpg
% pic07.jpg = 20261012_191636744.jpg
% pic08.jpg = 20261012_191643146.jpg
% pic09.jpg = 20261012_191649270.jpg
% pic10.jpg = 20261012_191656596.jpg
% pic11.jpg = 20261012_191702656.jpg
% Standalone XeLaTeX document, or input into an existing XePersian document.
\ifdefined\settextfont
  \def\HTNStandalone{0}
\else
  \def\HTNStandalone{1}
  \documentclass[12pt,a4paper]{article}
  \usepackage[margin=2cm]{geometry}
  \usepackage{amsmath,amssymb,array}
  \usepackage{xepersian}
  \IfFontExistsTF{Amiri}{\settextfont{Amiri}}{\settextfont{FreeSerif}}
  \setlatintextfont{FreeSerif}
\fi
\providecommand{\en}[1]{\lr{#1}}
\newcommand{\HTNarrow}{\en{$\longleftarrow$}}
\newcommand{\HTNcut}[1]{\textit{[متن لبهٔ تصویر: #1]}}
\newcommand{\HTNpage}[1]{\clearpage
  \section*{موضوع: براونوالد ـ ۲۶ ـ \en{HTN}}
  {\footnotesize تاریخ: \en{/ /}\quad صفحهٔ #1\par}}
\ifnum\HTNStandalone=1 \begin{document}\fi

\HTNpage{۱}
\subsection*{مقایسهٔ دو گایدلاین}
\begin{center}
\renewcommand{\arraystretch}{1.5}
\begin{tabular}{|r|c|c|}
\hline
 & آمریکا & اروپا \\\hline
مطب & \en{$S\geq130$}، \en{$D\geq80$} & \en{$S\geq140$}، \en{$D\geq90$} \\\hline
میانگین روز & \en{$S\geq130$}، \en{$D\geq80$} & \en{$S\geq135$}، \en{$D\geq85$} \\\hline
میانگین شب & \en{$S\geq110$}، \en{$D\geq65$} & \en{$S\geq120$}، \en{$D\geq70$} \\\hline
میانگین ۲۴ ساعت & \en{$S\geq125$}، \en{$D\geq75$} & \en{$S\geq130$}، \en{$D\geq80$} \\\hline
میانگین خانه & \en{$S\geq130$}، \en{$D\geq80$} & \en{$S\geq135$}، \en{$D\geq85$} \\\hline
هدف درمان & \en{$\boxed{<130/80}$} & \en{$S<140$} \\\hline
\end{tabular}
\end{center}
\subsection*{چگونه در خانه اندازه‌گیری کند؟}
\underline{۲ بار صبح}، \underline{۲ بار شب}، به مدت یک هفته، ۲ روز اول \HTNcut{«حذ…»؛ ادامهٔ واژه بیرون کادر است}.
\subsection*{اپروچ تشخیصی به پرایمری هایپرتنشن}
\begin{enumerate}
\item تأیید تشخیص، اندازه‌گیری مکرر فشار خون.
\item ارزیابی علل قابل رفع \en{HTN}.
\end{enumerate}

\HTNpage{۲}
\begin{enumerate}
\setcounter{enumi}{2}
\item شناسایی بیماری‌های همراه با \en{HTN} و کوموربیدیتی‌ها.
\item آموزش فاکتورهای مرتبط، سبک زندگی به بیمار.
\item ارزیابی آسیب ارگان‌های انتهایی در زمینهٔ \en{HTN}.
\item انتخاب نوع درمان.
\end{enumerate}
\subsection*{علل قابل رفع فشار خون}
\begin{itemize}
\item آپنهٔ خواب (خروپف، آپنه، لتارژیک).
\item هایپر / هایپوتیروئیدی.
\item هایپرپاراتیروئیدی (علائم هایپرکلسمی).
\item سندرم کوشینگ (شواهد کورتیزول \en{$\uparrow$}).
\item فئوکروموسایتوم، پاراگانگلیوم \HTNarrow{} علائم کاتکول‌آمین بالا.
\item آکرومگالی.
\end{itemize}
وقتی یک دست برای اندازه‌گیری فشار انتخاب شد باید دیگر از آن دست فشار گرفت.
\subsection*{آزمایش‌های اولیه}
\begin{latin}
\noindent CBC, BS, Ca, Na, K, BUN, Cr\\
Cholestrol, HDL, LDL, TG, U/A\\
\fbox{EKG}
\end{latin}

\HTNpage{۳}
شایع‌ترین آسیب ارگان در \en{HTN} \HTNarrow{} \en{\fbox{LVH}}.
\begin{itemize}
\item اکو به عنوان روتین در \en{HTN} مطرح نیست.
\end{itemize}
اگر بیمار در حال دریافت دیورتیک مصرف می‌کند برای ارزیابی سطح الکترولیت‌ها بهترین روش \underline{جمع‌آوری ادرار ۲۴ ساعته} است.

مهم است که بیمار \en{HTN} دچار هایپوکالمی (\en{$K\downarrow$}) نشود چون خود آن \en{BP} را بالا می‌برد.
\subsection*{پروفایل رنین}
توصیه می‌شود قبل از درمان اندازه‌گیری شود اگر به \HTNcut{ادامهٔ عبارت و یادداشت حاشیه‌ای در لبهٔ چپ بریده شده؛ چند حرف ناقص از یادداشت حاشیه‌ای قابل مشاهده است}.
\begin{enumerate}
\item \en{high plasma renin activity}:
\en{$>0.65\,\mathrm{ng/ml/h}$}، \en{$>6.5\,\mathrm{ng/ml/h}$}.
\par در این بیماران وازوکانستریکشن به علت فعالیت بیش از حد \en{RAAS} دارند.
\item \en{Suppresed plasma renin activity}:
\en{$<0.65\,\mathrm{ng/ml/h}$}.
\par مکانیسم فشار، \en{Volume Overload} است.
\end{enumerate}
\begin{enumerate}
\item \en{ACEI, ARB, ARNI, BB}.
\item \en{MRA, CCB, $\alpha B$}.
\end{enumerate}

\HTNpage{۴}
معمولاً با شرح حال و معاینه می‌فهمیم نیاز به اندازه‌گیری پروفایل رنین وجود دارد یا خیر.
\subsection*{فشار خون ثانویه}
\begin{itemize}
\item باید در همهٔ بیماران با \en{HTN} در نظر گرفته شود.
\item در بیمار با سابقهٔ \en{$\boxed{K\downarrow}$} باید حتماً در نظر گرفته شود.
\item معمولاً در زمینهٔ اختلالات اندوکراین هستند.
\item شایع‌ترین علت: \underline{پرایمری هایپرآلدو} (۲۵\%).
\item علل غیر اندوکراین: کوآرکتاسیون آئورت، رنوواسکولار دیزیز.
\end{itemize}
اگر در \underline{۲ نوبت} فشار به میزان \en{\fbox{20/10}} بالاتر از \en{130/80} بود در همان قدم اول، هم \en{lifestyle} هم دارو می‌دهیم.

اگر \en{\boxed{$eGFR<60$}} باشد حداقل ۲ دارو برای کنترل \en{BP} نیاز است.
\begin{itemize}
\item باید حتماً رژیم کم سدیم داشته باشد: \underline{کمتر از ۱٫۸ گرم} در ۲۴ ساعت.
\item مصرف \en{NSAID} باید در \underline{حداقل} خود باشد.
\item کاهش وزن به میزان \en{10 lb}.
\end{itemize}

\HTNpage{۵}
\en{NSAID} می‌تواند فشار را تا \en{6 mmHg} بالا ببرد.

اگر بیمار مجبور است \en{NSAID} مصرف کند بهتر است \en{\fbox{CCB}} و در درجهٔ بعد دیورتیک هم بگیرد تا اثر آن را مهار کند.
\subsection*{درمان در افراد با سن \en{$>65$}}
\begin{itemize}
\item درمان انتخابی اول در این بیماران \en{\fbox{CCB}} و دیورتیک \underline{تیازیدی} است.
\item دلیل: سطح رنین پایین در این سن؛ معمولاً مهارکننده‌های \en{RAAS} در این سن خیلی مؤثر نیستند.
\item در غیاب کوموربیدیتی \en{CV} نیاز به شروع \underline{\en{BB} نیست}.
\item اغلب بیماران ممکن است با \underline{یک دارو} کارشان راه بیفتد ولی در صورت \en{20/10} بالاتر از هدف، با \en{\fbox{2}} دارو شروع کنیم.
\item در سن بالا \HTNarrow{} \en{\fbox{CCB}} + تیازید (هیدروکلروتیازید / کلرتالیدون / \HTNcut{ادامهٔ نام دارو بریده شده است}).
\end{itemize}
عارضهٔ تیازید \HTNarrow{} \en{$K\downarrow$}، \en{$Mg\downarrow$}، \en{$Na\downarrow$}.

اگر بیمار سابقهٔ \underline{هایپوناترمی} (\en{$Na\downarrow$}) دارد یا در آزمایش‌های اولیه اختلال الکترولیتی دیده شود \HTNarrow{} بهتر است تیازید \underline{شروع نشود}.

\HTNpage{۶}
در طول درمان با تیازید، \underline{سطح پتاسیم} باید چک شود و در صورت افت، \underline{مکمل خوراکی} داده شود.
\begin{itemize}
\item در افراد مسن هدف درمان \underline{\en{$SBP<120$}}.
\item بجز افرادی که سفتی عروق دارند (\en{$PP>90\,\mathrm{mm}$}، فشار نبضی).
\item هدف \HTNarrow{} \underline{\en{$SBP<140$}}.
\end{itemize}
عارضهٔ \en{CCB} \HTNarrow{} ادم مچ پا، سردرد، هایپوتنشن وضعیتی.

در بیماران با \en{\fbox{LV dysfunction}} نباید \en{CCB} \underline{نسل اول} بدهیم.
\begin{itemize}
\item نیفدیپین، وراپامیل، دیلتیازم.
\end{itemize}
در بیماران با \underline{اختلال هدایت قلبی} نباید \en{Non-DHP CCB} داد چون هوا \HTNcut{خوانش واژهٔ «هوا» نامطمئن است} بلاک را بالا می‌برد.

بیماران بالای ۷۰ سال چون کم آب مصرف می‌کنند مستعد آسیب کلیه هستند و شروع مهارکنندهٔ \en{RAAS} دچار آسیب کلیه شوند.
\begin{itemize}
\item توصیه: افزایش مصرف آب.
\end{itemize}
\noindent $\star$ \en{BB} به عنوان مونوتراپی جایگاهی در \en{HTN} ندارد.

\HTNpage{۷}
اگر بخواهیم در بیمار مسن فشار خون به علت کوموربیدیتی‌های \en{CV}، \en{BB} بدهیم \HTNarrow{} \underline{\en{Nebivolol}} (سلکتیو \en{$\beta_1$}).
\begin{itemize}
\item عوارض افسردگی، \en{$BS\uparrow$}، دیس‌لیپیدمی و اختلال جنسی ندارد.
\end{itemize}
دیورتیک‌های پتاسیم اسپیرینگ (در منبع: \en{MRA} خط خورده) فقط در \underline{\en{$eGFR>45$}} قابل تجویز \HTNarrow{} خطر هایپرکالمی.
% A crossed-out abbreviation resembling MRA appears before «فقط».

\en{MRA} \HTNarrow{} در بیمار \en{HTN + \fbox{HF}} یا پرایمری هایپرآلدو.
\begin{itemize}
\item عوارض (اسپیرونولاکتون): ژنیکوماستی، هیرسوتیسم.
\end{itemize}
مهارکننده‌های آلفا (کلونیدین) و وازودیلاتورها غیر اختصاصی (ماینوکسیدیل) نباید به عنوان خط اول درمان \en{HTN} استفاده شوند در بیماران مسن.
\subsection*{تجویز \en{HTN} در بیمار \en{CKD}}
\begin{itemize}
\item دو هدف \en{$<130/80$} و \en{120/80} \underline{برای کلیه تفاوت ندارند} ولی \underline{برای \en{CV} چرا}.
\end{itemize}
ریسک فاکتورهای هایپرکالمی (\en{$K\uparrow$}) در بیمار \en{CKD}:
\begin{enumerate}
\item \underline{\en{$eGFR<45$}}.
\item \en{$K>4.5$}.
\item \en{$BMI<25$}.
\end{enumerate}

\HTNpage{۸}
\subsection*{سبک زندگی در کنترل \en{BP} در \en{CKD}}
\begin{itemize}
\item مصرف روزانهٔ نمک کمتر از \en{2.4 gr}.
\item کاهش مصرف الکل.
\item ورزش ایروبیک \en{$\checkmark$}، ورزش ایزومتریک \en{$\times$}.
\end{itemize}
اگر \en{$eGFR<90$} باشد شروع درمان با \underline{یک دارو} بهتر است.

در بیمار با \underline{نفروپاتی دیابتی} باید با مهارکنندهٔ \en{RAAS} درمان را شروع کنیم.
\begin{itemize}
\item اگر نفروپاتی نباشد یا آلبومینوری \en{$<300\,\mathrm{mg}$} روزانه نداشته باشد مهارکنندهٔ \en{RAAS} \underline{برتری ندارد}.
\end{itemize}
در بیمار نورموتنسیو \en{$\pm$} میکروآلبومینوری، تجویز مهارکنندهٔ \en{RAAS} نقشی در پیشگیری از افت عملکرد کلیه \underline{نمی‌کند}.
\subsection*{تعریف \en{HTN} مقاوم}
شکست در رساندن \en{BP} به \underline{\en{$<140/90$}} در بیمار تحت درمان با حداکثر دوز ۳ داروی همزمان فشار خون که حداقل یکی از آن‌ها دیورتیک است.
\begin{itemize}
\item ۸–۱۲\% بیماران \en{HTN}، مقاوم هستند.
\end{itemize}

\HTNpage{۹}
\begin{itemize}
\item درمان اضافه کردن داروی چهارم است.
\item \en{CCB + RAAS Inh} + دیورتیک.
\item \en{MRA} (اسپیرونولاکتون).
\end{itemize}
\subsection*{اورژانس فشار خون}
ترکیب فشار بالا (حد مشخصی ندارد) + علائم و نشانه‌های
\begin{latin}
\noindent Acute, Ongoing End organ damage
\end{latin}
هایپرتنسیو ارجنسی \HTNarrow{} \en{$SBP>180$} بدون آسیب حاد ارگان‌های انتهایی.
\begin{itemize}
\item \underline{درمان خوراکی} با یک یا چند دارو و پیگیری.
\end{itemize}
\subsection*{ارزیابی بیمار اورژانس فشار خون}
\begin{itemize}
\item معاینهٔ حجم، \en{MS}، قلبی، ریوی، نورولوژیک.
\item بروئی شکمی.
\end{itemize}
\subsection*{اورژانس فشار خون در شرایط مختلف}
\begin{itemize}
\item دایسکشن حاد آئورت \HTNarrow{} ظرف \en{\fbox{20 min}}، \en{$\boxed{SBP<120}$}.
\item \en{BB} (اثری مشابه وراپامیل / دیلتیازم دارند) + وازودیلاتور.
\end{itemize}

\HTNpage{۱۰}
\subsection*{استروک ایسکمیک حاد}
\begin{itemize}
\item اگر بیمار کاندید ترومبولیتیک است و فشار \underline{\en{$>180/110$}} است، فشار سریعاً پایین آورده شود.
\item بیمار استروک را صرفاً زمانی فشار را پایین می‌آوریم که بالای \en{180/110} باشد که آن هم روی \en{outcome} استروک تأثیری ندارد.
\end{itemize}
\subsection*{سایر انواع اورژانس‌ها}
\begin{itemize}
\item \underline{۱۰–۱۵\% طی ساعت اول}.
\item ۱۰–۲۰\% طی یک ساعت بعدی.
\item \underline{جمعاً ۲۵\% طی ۲۴ ساعت اول}.
\item انسفالوپاتی هایپرتنسیو: کنترل فشار به سرعت خوب می‌شود.
\end{itemize}
\noindent $\star$ \en{ACS} یا ادم ریوی:
\begin{itemize}
\item نیتروگلیسیرین، \en{Clevidipine}، نیکاردیپین، نیتروپروساید.
\item معمولاً \underline{\en{ACEI}} هم به خاطر بیماری دریافت خواهد کرد \HTNcut{انتهای خط تا لبهٔ چپ می‌رود؛ خوانش «خواهد کرد» نامطمئن است}.
\end{itemize}
\en{AKI} \HTNarrow{} \underline{فنولدوپام} ارجح؛ نیکاردیپین و نیتروپروساید.

افزایش کاتکول‌آمین (مسمومیت، فئو، \en{MAOI}):
\begin{itemize}
\item \en{$\boxed{\alpha B}$} (فنتولامین) + \en{BB}.
\end{itemize}
قطع کلونیدین ناگهانی \HTNarrow{} یک دوز کلونیدین.

حین بارداری \HTNarrow{} منیزیم سولفات، متیل‌دوپا، \underline{هیدرالازین}.
\begin{itemize}
\item خط دوم: لابتالول خوراکی، نیفدیپین.
\end{itemize}

\HTNpage{۱۱}
داروی خوراکی باید ظرف \underline{۸–۲۴ ساعت} پس از شروع درمان تزریقی آغاز گردد.

دوز نیتروپروساید \en{$[0.5\,\mu\mathrm{g/kg/min}]$}.

\en{HTN Ergency} \HTNarrow{} کلونیدین، کاپتوپریل، لابتالول، آملودیپین.
\begin{itemize}
\item ترجیحاً نیفدیپین ندهیم.
\end{itemize}
فنولدوپام \HTNarrow{} آگونیست دوپامین ـ ۱، عملکرد کلیه را بهبود می‌بخشد.

\en{Clevidipine} \HTNarrow{} دفع آن تحت تأثیر عملکرد کبد و کلیه نیست.
\begin{itemize}
\item $\star$ به دلیل ترکیبات همراه ریسک آلرژی بالا.
\item در \en{ACS} چون رفلکس تاکی دارد کاربردی \underline{ندارد}.
\end{itemize}
نیمودیپین \HTNarrow{} فقط در \en{SAH} جایگاه دارد.

% Printed subject/date labels and handwritten page numbers are preserved.
% Notebook emblems contain no text. Clipped/unreadable source text is marked.
\ifnum\HTNStandalone=1
  \let\HTNEnd\enddocument
\else
  \let\HTNEnd\relax
\fi
\HTNEnd
